\SixthArticle{33}
Just as in the preceding article we reduced $\Psi\left(-\tfrac12\right)$ to $\Psi0$ and a logarithm, so in general we shall reduce $\Psi\left(-\frac mn\right)$, with $m,n$ denoting integers of which $m$ is the smaller, to $\Psi0$ and logarithms. Let us set $\frac{2\pi}{n}=\omega$, and let $\varphi$ equal one of the angles $\omega,2\omega,3\omega\ .\ .\ .\ .\ (n-1)\omega$; hence $1=\cos n\varphi=\cos2n\varphi=\cos3n\varphi$ etc., $\cos\varphi=\cos(n+1)\varphi=\cos(n+2)\varphi$ etc., $\cos2\varphi=\cos(n+2)\varphi$ etc., and also $\cos\varphi+\cos2\varphi+\cos3\varphi+\text{etc.}+\cos(n-1)\varphi+1=0$. We therefore have
\[
\begin{aligned}
&\cos\varphi\cdot\Psi\frac{1-n}{n}=-n\cos\varphi+\cos\varphi\cdot\log2\\[4pt]
&-\frac{n}{n+1}\cos(n+1)\varphi+\cos\varphi\cdot\log\tfrac32-\text{etc.}
\end{aligned}
\]
\[
\begin{aligned}
&\cos2\varphi\cdot\Psi\frac{2-n}{n}=-\frac n2\cos2\varphi+\cos2\varphi\cdot\log2\\[4pt]
&-\frac{n}{n+2}\cos(n+2)\varphi+\cos2\varphi\cdot\log\tfrac32-\text{etc.}
\end{aligned}
\]
\[
\begin{aligned}
&\cos3\varphi\cdot\Psi\frac{3-n}{n}=-\frac n3\cos3\varphi+\cos3\varphi\cdot\log2\\[4pt]
&-\frac{n}{n+3}\cos(n+3)\varphi+\cos3\varphi\cdot\log\tfrac32-\text{etc.}
\end{aligned}
\]
etc. up to
\SourceLeaf{165}
\[
\begin{aligned}
&\cos(n-1)\varphi\cdot\Psi\left(-\frac1n\right)
\\[4pt]
&=-\frac{n}{n-1}\cos(n-1)\varphi+\cos(n-1)\varphi\cdot\log2\\[4pt]
&-\frac{n}{2n-1}\cos(2n-1)\varphi
+\cos(n-1)\varphi\cdot\log\tfrac32-\text{etc.}
\end{aligned}
\]
\[
\Psi0=-\frac nn\cos n\varphi+\log2-\frac{n}{2n}\cos2n\varphi+\log\tfrac32-\text{etc.}
\]
and by \textit{summation}
\[
\begin{aligned}
&\cos\varphi\cdot\Psi\frac{1-n}{n}+\cos2\varphi\cdot\Psi\frac{2-n}{n}\\[4pt]
&+\cos3\varphi\cdot\Psi\frac{3-n}{n}+\text{etc.}\\[4pt]
&+\cos(n-1)\varphi\cdot\Psi\left(-\frac1n\right)+\Psi0
\\[4pt]
&=-n(\cos\varphi+\tfrac12\cos2\varphi+\tfrac13\cos3\varphi+\tfrac14\cos4\varphi+\text{etc. to infinity})
\end{aligned}
\]
But in general, for a value of $x$ no greater than unity,
\[
\begin{aligned}
&\log(1-2x\cos\varphi+xx)\\[4pt]
&=-2(x\cos\varphi+\tfrac12xx\cos2\varphi+\tfrac13x^3\cos3\varphi+\text{etc.})
\end{aligned}
\]
This series follows easily from the expansion of $\log(1-rx)+\log\left(1-\frac xr\right)$, where $r$ denotes the quantity $\cos\varphi+\BareRoot{-1}\cdot\sin\varphi$. Thus the preceding equation becomes
\[
\begin{aligned}
&\cos\varphi\cdot\Psi\frac{1-n}{n}+\cos2\varphi\cdot\Psi\frac{2-n}{n}\\[4pt]
&+\cos3\varphi\cdot\Psi\frac{3-n}{n}+\text{etc.}\\[4pt]
&+\cos(n-1)\varphi\cdot\Psi\left(-\frac1n\right)
\\[4pt]
&=-\Psi0+\tfrac12n\log(2-2\cos\varphi)
\end{aligned}
\tag*{[72]}\]
In this equation let $\varphi=\omega$, $\varphi=2\omega$, $\varphi=3\omega$ etc. up to $\varphi=(n-1)\omega$ be set successively; let these equations be multiplied in their order by $\cos m\omega$, $\cos2m\omega$, $\cos3m\omega$ etc. up to $\cos(n-1)m\omega$, and let equation 71 be added to the sum of the products:
\[
\begin{aligned}
&\Psi\frac{1-n}{n}+\Psi\frac{2-n}{n}+\Psi\frac{3-n}{n}+\text{etc.}+\Psi\left(-\frac1n\right)\\[4pt]
&=(n-1)\Psi0-n\log n
\end{aligned}
\]
Now, on considering that
\[
\begin{aligned}
&1+\cos m\omega\cdot\cos k\omega+\cos2m\omega\cdot\cos2k\omega\\[4pt]
&+\cos3m\omega\cdot\cos3k\omega
+\text{etc.}+\cos(n-1)m\omega\cdot\cos(n-1)k\omega=0
\end{aligned}
\]
where $k$ denotes any of the numbers $1,2,3\ .\ .\ .\ .\ (n-1)$ except the two $m$ and $n-m$, for which that sum is $=\tfrac12n$, it will be clear that summing those equations and dividing by $\frac n2$ gives
\[
\begin{aligned}
&\Psi\left(-\frac mn\right)+\Psi\left(-\frac{n-m}{m}\right)=
\\[4pt]
&2\Psi0-2\log n+\cos m\omega\cdot\log(2-2\cos\omega)\\[4pt]
&+\cos2m\omega\cdot\log(2-2\cos2\omega)
\\[4pt]
&+\cos3m\omega\cdot\log(2-2\cos3\omega)+\text{etc.}\\[4pt]
&+\cos(n-1)m\omega\cdot\log(2-2\cos(n-1)\omega)
\end{aligned}
\tag*{[73]}\]
Evidently the last term of this equation is $=\cos m\omega\cdot\log(2-2\cos\omega)$, the\SourceLeaf{166} next to last $=\cos 2m\omega.\log(2-2\cos 2\omega)$ etc., so that the terms are always equal in pairs, except, if $n$ is even, for the single term $\cos\frac{n}{2}.m\omega\log(2-2\cos\frac{n}{2}\omega)$, which is $=+2\log 2$ for even $m$, or $=-2\log 2$ for odd $m$. Combining with equation 73 the equation
\[
\Psi\left(-\frac{m}{n}\right)-\Psi\left(-\frac{n-m}{n}\right)=\pi\operatorname{cotang}\frac{m}{n}\pi
\]
we have, for an odd value of $n$, provided $m$ is a positive integer smaller \mbox{than $n$},
\[
\begin{aligned}
&{}\text{[74]}\qquad \Psi\left(-\frac{m}{n}\right)
\\[4pt]
&=\Psi 0+\tfrac12\pi\operatorname{cotang}\frac{m\pi}{n}-\log n\\[4pt]
&+\cos\frac{2m\pi}{n}.\log\left(2-2\cos\frac{2\pi}{n}\right)
\\[4pt]
&+\cos\frac{4m\pi}{n}.\log\left(2-2\cos\frac{4\pi}{n}\right)\\[4pt]
&+\cos\frac{6m\pi}{n}.\log\left(2-2\cos\frac{6\pi}{n}\right)+\text{etc.}
\\[4pt]
&+\cos\frac{(n-1)m\pi}{n}.\log\left(2-2\cos\frac{(n-1)\pi}{n}\right)
\end{aligned}
\]
For an even value of $n$, however,
\[
\begin{aligned}
&{}\text{[75]}\qquad \Psi\left(-\frac{m}{n}\right)
\\[4pt]
&=\Psi 0+\tfrac12\pi\operatorname{cotang}\frac{m\pi}{n}-\log n\\[4pt]
&+\cos\frac{2m\pi}{n}\log\left(2-2\cos\frac{2\pi}{n}\right)
\\[4pt]
&+\cos\frac{4m\pi}{n}\log\left(2-2\cos\frac{4\pi}{n}\right)+\text{etc.}\\[4pt]
&+\cos\frac{(n-2)m\pi}{n}\log\left(2-2\cos\frac{(n-2)\pi}{n}\right)
\pm\log 2
\end{aligned}
\]
where the upper sign holds for even $m$, the lower for odd. Thus, for example,
\[
\begin{aligned}
&\Psi\left(-\tfrac14\right)=\Psi 0+\tfrac12\pi-3\log 2, 
\\[4pt]
&\Psi\left(-\tfrac34\right)=\Psi 0-\tfrac12\pi-3\log 2
\\[4pt]
&\Psi\left(-\tfrac13\right)=\Psi 0+\tfrac12\pi\BareRoot{\tfrac13}-\tfrac32\log 3, 
\\[4pt]
&\Psi\left(-\tfrac23\right)=\Psi 0-\tfrac12\pi\BareRoot{\tfrac13}-\tfrac22\log 3
\end{aligned}
\]
Moreover, combining these equations with equation 64 shows of itself that, in general, $\Psi z$ can be determined for \emph{any rational value} of $z$, positive or negative, by $\Psi 0$ and logarithms; this theorem is certainly most remarkable.

\SixthArticle{34}
Since, by Article 28, $\Pi\lambda$ is the value of the integral $\int y^\lambda e^{-y}\mathrm{d}y$ from $y=0$ to $y=\infty$, provided $\lambda+1$ is a positive quantity, differentiation with respect to $\lambda$ gives
\[
\frac{\mathrm{d}\Pi\lambda}{\mathrm{d}\lambda}
=\frac{\mathrm{d}\int y^\lambda e^{-y}\mathrm{d}y}{\mathrm{d}\lambda}
=\int y^\lambda e^{-y}\log y\,\mathrm{d}y
\]
or
\[
\begin{aligned}
&[76]\qquad \Pi\lambda.\Psi\lambda=\int y^\lambda e^{-y}\log y.\mathrm{d}y,\\[4pt]
&\quad\text{from }y=0\text{ to }y=\infty
\end{aligned}
\]
\par\SourceLeaf{167}
{\raggedright\parindent=1.2em More generally, on setting $y=z^\alpha$, $\alpha\lambda+\alpha-1=\beta$, the value of the integral $\int z^\beta e^{-z^\alpha}\log z\,\mathrm{d}z$ from $z=0$ to $z=\infty$ is\par}
\[
\begin{aligned}
&=\frac{1}{\alpha\alpha}\Pi\left(\frac{\beta+1}{\alpha}-1\right).\Psi\left(\frac{\beta+1}{\alpha}-1\right)
\\[4pt]
&=\frac{1}{\alpha(\beta+1)}\Pi\frac{\beta+1}{\alpha}.\Psi\frac{\beta+1}{\alpha}
-\frac{1}{(\beta+1)^2}\Pi\frac{\beta+1}{\alpha}
\end{aligned}
\]
provided $\beta+1$ and $\alpha$ are both positive quantities; or it equals the same quantity with the opposite sign if both $\beta+1$ and $\alpha$ are negative.

